\documentclass[10pt]{article}
\usepackage{amsmath,amssymb,amsfonts}
\usepackage{geometry}
\usepackage{xcolor}
\usepackage{tcolorbox}
\tcbuselibrary{skins,breakable}
\usepackage{tikz}
\usepackage{fancyhdr}
\usepackage{enumitem}
\usepackage{newtxtext,newtxmath}

\geometry{
    paperwidth=8.5in,
    paperheight=11in,
    top=0.55in,
    bottom=0.55in,
    left=0.75in,
    right=0.75in,
    headheight=14pt,
    headsep=12pt,
    footskip=20pt
}

% Định nghĩa màu chuẩn sách Stewart Calculus
\definecolor{stewartcyan}{RGB}{0, 136, 170}
\definecolor{stewartred}{RGB}{204, 34, 41}
\definecolor{projectbg}{RGB}{236, 243, 236}

% Định nghĩa các hàm toán học
\DeclareMathOperator{\sech}{sech}
\DeclareMathOperator{\csch}{csch}
\DeclareMathOperator{\coth}{coth}

% Cấu hình khung định nghĩa
\newtcolorbox{definitionbox}[1][]{
    enhanced,
    colback=white,
    colframe=stewartred,
    boxrule=1.1pt,
    arc=0mm,
    left=14pt,
    right=14pt,
    top=8pt,
    bottom=8pt,
    boxsep=0pt,
    #1
}

% Thiết lập đầu trang và chân trang
\pagestyle{fancy}
\fancyhf{}
\renewcommand{\headrulewidth}{0pt}
\fancyhead[R]{\textsf{\textbf{MỤC 3.11}\quad Các hàm Hyperbol}\qquad\textbf{\textsf{261}}}
\fancyfoot[C]{%
    \fontsize{6.5}{8}\selectfont\color{black!60}%
    Bản quyền 2021 Cengage Learning. Bảo lưu mọi quyền. Không được sao chép, quét hoặc nhân bản toàn bộ hay một phần. WCN 02-200-203\\[2pt]
    Đánh giá biên tập cho thấy mọi nội dung bị loại bỏ không ảnh hưởng đáng kể đến trải nghiệm học tập tổng thể. Cengage Learning bảo lưu quyền xóa nội dung bổ sung nếu các hạn chế về quyền sau này yêu cầu.
}

\begin{document}
\thispagestyle{fancy}

% PHẦN DỰ ÁN ỨNG DỤNG TIẾP DIỄN (PROJECT CONTINUATION)
\begin{tcolorbox}[
    enhanced,
    colback=projectbg,
    colframe=projectbg,
    arc=0mm,
    boxrule=0pt,
    left=0pt,
    right=0pt,
    top=8pt,
    bottom=10pt,
    boxsep=0pt
]
\noindent
\begin{minipage}[t]{120pt}
    \null % Cột lề trống chuẩn layout hai cột của Stewart
\end{minipage}%
\hspace{16pt}%
\begin{minipage}[t]{368pt}
    \setlength{\abovedisplayskip}{3pt}
    \setlength{\belowdisplayskip}{3pt}
    \fontsize{9.5}{13.5}\selectfont
    \noindent Hãy chứng minh rằng hàm bậc hai thỏa mãn các điều kiện (i), (ii) và (iii) là
    \[
    P(x) = f(a) + f'(a)(x - a) + \tfrac{1}{2} f''(a)(x - a)^2
    \]

    \begin{enumerate}[label=\textbf{\arabic*.}, leftmargin=1.4em, itemsep=5pt, topsep=3pt, parsep=0pt, start=4]
        \item Tìm xấp xỉ bậc hai của hàm số $f(x) = \sqrt{x + 3}$ gần $a = 1$. Vẽ đồ thị của $f$, xấp xỉ bậc hai và xấp xỉ tuyến tính từ Ví dụ 3.10.2 trên cùng một màn hình. Bạn rút ra kết luận gì?

        \item Thay vì chỉ thỏa mãn với một xấp xỉ tuyến tính hoặc xấp xỉ bậc hai cho $f(x)$ gần $x = a$, chúng ta hãy thử tìm các xấp xỉ tốt hơn bằng các đa thức bậc cao hơn. Ta tìm một đa thức bậc $n$
        \[
        T_n(x) = c_0 + c_1(x - a) + c_2(x - a)^2 + c_3(x - a)^3 + \dots + c_n(x - a)^n
        \]
        sao cho $T_n$ và $n$ đạo hàm đầu tiên của nó có cùng giá trị tại $x = a$ với $f$ và $n$ đạo hàm đầu tiên của $f$. Bằng cách lấy đạo hàm liên tiếp và cho $x = a$, hãy chứng minh rằng các điều kiện này được thỏa mãn nếu $c_0 = f(a)$, $c_1 = f'(a)$, $c_2 = \tfrac{1}{2} f''(a)$, và tổng quát
        \[
        c_k = \frac{f^{(k)}(a)}{k!}
        \]
        trong đó $k! = 1 \cdot 2 \cdot 3 \cdot 4 \dotsm k$. Đa thức thu được
        \[
        T_n(x) = f(a) + f'(a)(x - a) + \frac{f''(a)}{2!}(x - a)^2 + \dots + \frac{f^{(n)}(a)}{n!}(x - a)^n
        \]
        được gọi là \textbf{đa thức Taylor bậc $\boldsymbol{n}$ của $\boldsymbol{f}$ tại $\boldsymbol{a}$}. (Chúng ta sẽ nghiên cứu các đa thức Taylor chi tiết hơn trong Chương 11.)

        \item Tìm đa thức Taylor bậc 8 tại $a = 0$ cho hàm số $f(x) = \cos x$. Vẽ đồ thị của $f$ cùng với các đa thức Taylor $T_2, T_4, T_6, T_8$ trong khung nhìn $[-5, 5] \times [-1{,}4;\, 1{,}4]$ và nhận xét về mức độ xấp xỉ tốt của chúng đối với $f$.
    \end{enumerate}
\end{minipage}
\end{tcolorbox}

\vspace{14pt}

% PHẦN TIÊU ĐỀ MỤC 3.11
\noindent
\begin{minipage}[b]{120pt}
    \raggedleft
    \begin{tikzpicture}[baseline={(0, 3pt)}]
        \foreach \y in {0pt, 3.2pt, 6.4pt, 9.6pt, 12.8pt} {
            \draw[stewartcyan, line width=0.85pt] (0pt, \y) -- (44pt, \y);
        }
    \end{tikzpicture}\hspace{7pt}%
    {\fontsize{22}{22}\selectfont\sffamily\bfseries\color{stewartcyan}3.11}
\end{minipage}%
\hspace{16pt}%
\begin{minipage}[b]{368pt}
    {\fontsize{20}{22}\selectfont\sffamily\bfseries Các hàm Hyperbol}
\end{minipage}

\vspace{12pt}

% PHẦN NỘI DUNG MỤC 3.11
\noindent
\begin{minipage}[t]{120pt}
    \null
\end{minipage}%
\hspace{16pt}%
\begin{minipage}[t]{368pt}
    \noindent{\color{stewartcyan}\rule{6.5pt}{6.5pt}}\hspace{6pt}{\fontsize{11.5}{14}\selectfont\sffamily\bfseries Các hàm Hyperbol và đạo hàm của chúng}\par
    \vspace{7pt}

    \fontsize{9.5}{13.5}\selectfont
    Một số tổ hợp nhất định của các hàm mũ $e^x$ và $e^{-x}$ xuất hiện thường xuyên trong toán học và các ứng dụng của nó đến mức chúng xứng đáng được đặt những tên gọi riêng. Về nhiều mặt, chúng tương tự như các hàm lượng giác, và chúng có cùng mối quan hệ với hypebol như mối quan hệ của các hàm lượng giác đối với đường tròn. Vì lý do này, chúng được gọi chung là \textbf{các hàm hyperbol} và được gọi riêng lẻ là \textbf{sin hyperbol}, \textbf{cosin hyperbol}, và tương tự như vậy.
    \vspace{10pt}

    \begin{definitionbox}
        {\sffamily\bfseries\color{stewartred}\fontsize{10.5}{13}\selectfont Định nghĩa các hàm Hyperbol}\par
        \vspace{4pt}
        \begin{alignat*}{2}
            \sinh x &= \frac{e^x - e^{-x}}{2} &\hspace{4.5em} \csch x &= \frac{1}{\sinh x} \\[8pt]
            \cosh x &= \frac{e^x + e^{-x}}{2} &\hspace{4.5em} \sech x &= \frac{1}{\cosh x} \\[8pt]
            \tanh x &= \frac{\sinh x}{\cosh x} &\hspace{4.5em} \coth x &= \frac{\cosh x}{\sinh x}
        \end{alignat*}
    \end{definitionbox}
\end{minipage}

\end{document}